\begin{examplebox}
	\textbf{\textcolor{CExample}{例 1（基础）}}
	\textbf{\textcolor{CExample}{题目：}}计算 $C(8,3)$ 和 $C(8,5)$，并说明它们的关系。\textbf{\textcolor{CExample}{解题步骤：}}
	\begin{description}[style=nextline, leftmargin=2.8em, labelsep=0.5em, itemsep=0.45em]
		\item[\textbf{第一步（代入公式）：}]
			\[
				C(8,3)=\frac{8!}{3!\,5!}, \quad C(8,5)=\frac{8!}{5!\,3!}
			\]
			\par\textbf{理由：}直接套用组合数公式。
		\item[\textbf{第二步（计算约分）：}]
			\[
				\begin{aligned}
					C(8,3) &= \frac{8\times7\times6}{3\times2\times1}\\ &= 56,
			\end{aligned}
				\qquad
				\begin{aligned}
					C(8,5) &= \frac{8\times7\times6}{3\times2\times1}\\ &= 56.
			\end{aligned}
			\]
			\par\textbf{理由：}展开阶乘后约分得到结果。
		\item[\textbf{第三步（结论验证）：}]
			$C(8,3)=C(8,5)$，验证了组合数的对称性质 $C(n,m)=C(n,n-m)$。
			\par\textbf{理由：}对称性得证。
	\end{description}
	\textbf{\textcolor{CExample}{关键结论：}}
	$C(8,3)=C(8,5)=56$，且 $\mathbf{C(n,m)=C(n,n-m)}$ 成立。

	\medskip
	\textbf{\textcolor{CExample}{例 2（稍变形）}}
	\textbf{\textcolor{CExample}{题目：}}已知 $C(n,2)=15$，求正整数 $n$。\textbf{\textcolor{CExample}{解题步骤：}}
	\begin{description}[style=nextline, leftmargin=2.8em, labelsep=0.5em, itemsep=0.45em]
		\item[\textbf{第一步（建立方程）：}]
			由 $C(n,2)=\dfrac{n(n-1)}{2}$ 且已知 $C(n,2)=15$，得方程 $\dfrac{n(n-1)}{2}=15$。
			\par\textbf{理由：}组合数公式导出 $C(n,2)=\frac{n(n-1)}{2}$。
		\item[\textbf{第二步（解二次数）：}]
			$n^2-n-30=0$，解得 $n=6$ 或 $n=-5$（舍去）。
			\par\textbf{理由：}$n$ 为正整数，取 $n=6$。
		\item[\textbf{第三步（验证结果）：}]
			$C(6,2)=\frac{6\times5}{2}=15$，符合条件。
			\par\textbf{理由：}代入验证。
	\end{description}
	\textbf{\textcolor{CExample}{关键结论：}}
	$n=6$，且使用公式 $C(n,2)=\frac{n(n-1)}{2}$ 可快速求解。

	\medskip
	\textbf{\textcolor{CExample}{例 3（综合应用）}}
	\textbf{\textcolor{CExample}{题目：}}盒中有 10 个球，其中 3 个红球，7 个蓝球。现从中任取 3 个球，求恰好取出 2 个红球的取法总数。\textbf{\textcolor{CExample}{解题步骤：}}
	\begin{description}[style=nextline, leftmargin=2.8em, labelsep=0.5em, itemsep=0.45em]
		\item[\textbf{第一步（总取法）：}]
			从 10 个球中取 3 个，总取法数 $C(10,3)$ 暂不计算。
			\par\textbf{理由：}本题只需求符合条件的方法数，不用求概率。
		\item[\textbf{第二步（分类计数）：}]
			取到 2 个红球：从 3 个红球中取 2 个有 $C(3,2)$ 种；同时从 7 个蓝球中取 1 个有 $C(7,1)$ 种。
			\par\textbf{理由：}分步计数（乘法原理）。
		\item[\textbf{第三步（计算数值）：}]
			\[
				\begin{aligned}
					C(3,2)\times C(7,1) &= \dfrac{3!}{2!\,1!}\times 7 \\ &= 3\times 7 \\ &= 21\text{ 种}.
			\end{aligned}
			\]
			\par\textbf{理由：}代入组合数公式并约分。
	\end{description}
	\textbf{\textcolor{CExample}{关键结论：}}
	恰好 2 个红球的取法数为 $21$，直接应用了组合数公式 $C(3,2)=\frac{3!}{2!1!}=3$ 和 $C(7,1)=7$。
\end{examplebox}
